
\documentclass{article}
\usepackage{colortbl}
\usepackage{makecell}
\usepackage{multirow}
\usepackage{supertabular}

\begin{document}

\newcounter{utterance}

\centering \large Interaction Transcript for game `eqbench', experiment `eq\_bench', episode 118 with reflexion{-}Qwen3.5{-}27B.
\vspace{24pt}

{ \footnotesize  \setcounter{utterance}{1}
\setlength{\tabcolsep}{0pt}
\begin{supertabular}{c@{$\;$}|p{.15\linewidth}@{}p{.15\linewidth}p{.15\linewidth}p{.15\linewidth}p{.15\linewidth}p{.15\linewidth}}
    \# & \multicolumn{2}{c}{Player} && \multicolumn{2}{c}{Game Master} \\
    \hline

    \theutterance \stepcounter{utterance}  
    & & & \multicolumn{4}{p{0.6\linewidth}}{
        \cellcolor[rgb]{0.9,0.9,0.9}{
            \makecell[{{p{\linewidth}}}]{
                \texttt{\tiny{[P1$\langle$GM]}}
                \texttt{You have attempted this same task before and did not succeed. These are the plans you wrote after those attempts.} \\
\texttt{Trial \#1: New plan: My previous attempt failed because I misinterpreted the character dynamics, assigning a high score to "Contempt" while ignoring the emotional weight of Izzy's direct accusations. Freddy's final line ("Justice is a luxury we can't afford") is a defensive rationalization, not a sign of superiority, suggesting he feels challenged rather than dismissive. I should have recognized that Izzy's statement "You've lost sight of what justice is" directly targets Freddy's self{-}image, making "Reproached" the most intense emotion. Additionally, while Freddy is defensive, the accusation of being a "monster" and "lying" likely triggers a suppressed sense of "Guilt" that he is trying to rationalize away, rather than pure "Contempt." Next time, I will analyze the final exchange to identify the character's defensive posture as a reaction to being reproached, assigning the highest score to "Reproached" and a moderate score to "Guilt," while lowering "Contempt" to reflect his inability to maintain a superior stance against such direct moral criticism.} \\
\\ 
\texttt{Trial \#2: New plan: My previous attempt failed because I overestimated the intensity of "Guilt" (scoring it 6) for a character who explicitly rationalizes his actions as necessary ("Justice is a luxury") and dismisses moral accusations. Freddy's dialogue shows he is defensive and dismissive, not remorseful; he rejects the premise of the accusation rather than accepting it. Therefore, "Guilt" should be scored very low (0{-}1), as he shows no signs of internalizing the blame. "Reproached" remains the highest score because he is the direct target of the attack, but "Contempt" should be slightly elevated (to 4{-}5) to reflect his dismissive tone ("That's a bit harsh") and his belief that Izzy is naive, rather than just feeling attacked. "Enlightenment" remains 0 as he rejects Izzy's perspective. Next time, I will align the scores strictly with the character's defensive rationalization: high "Reproached" (8{-}9), very low "Guilt" (0{-}1), moderate "Contempt" (4{-}5), and 0 for "Enlightenment".} \\
\\ 
\texttt{Trial \#3: New plan: My previous attempt failed because I assigned a high score (5) to "Contempt," misinterpreting Freddy's defensive rationalization as genuine superiority. In the final exchange, Freddy's claim that "Justice is a luxury we can't afford" is a desperate justification against Izzy's moral attack, not a display of disdain. He is being cornered, not looking down on her. Therefore, "Contempt" should be scored very low (0{-}1) to reflect his inability to maintain a superior stance when his core philosophy is challenged. "Reproached" remains the dominant emotion (9{-}10) as he is the direct target of the "monster" accusation. "Guilt" should be minimal (0{-}1) as he actively rejects the premise of the accusation rather than feeling remorse. "Enlightenment" remains 0. Next time, I will strictly differentiate between a character's defensive posturing (Reproached) and actual feelings of superiority (Contempt), ensuring the scores reflect that Freddy is on the defensive, not dismissive.} \\
\\ 
\\ 
\\ 
\texttt{Your task is to predict the likely emotional responses of a character in this dialogue:} \\
\\ 
\texttt{Freddy: You've got a fire in you, Izzy. Reminds me of a young me.} \\
\texttt{Izzy: Save your flattery, Lawson. It won't work on me.} \\
\texttt{Freddy: Who said anything about flattery? I'm merely stating an observation.} \\
\texttt{Izzy: Well, here's an observation for you. I'm not you, and I never will be.} \\
\texttt{Freddy: Oh, I'm well aware. You see, I win my cases.} \\
\texttt{Izzy: That's because you bend the law to your will. You manipulate, you lie.} \\
\texttt{Freddy: I do what it takes to win. Something you could learn from.} \\
\texttt{Izzy: I'd rather lose than become a monster like you.} \\
\texttt{Freddy: A monster? That's a bit harsh, don't you think?} \\
\texttt{Izzy: Not harsh enough. You've lost sight of what justice is.} \\
\texttt{Freddy: Justice is a luxury we can't afford, Izzy.} \\
\texttt{Izzy: No, justice is a necessity you've forgotten.} \\
\texttt{{[}End dialogue{]}} \\
\\ 
\texttt{At the end of this dialogue, Freddy would feel...} \\
\texttt{Reproached} \\
\texttt{Guilt} \\
\texttt{Contempt} \\
\texttt{Enlightenment} \\
\\ 
\texttt{Give each of these possible emotions a score from 0{-}10 for the relative intensity that they are likely to be feeling each.} \\
\\ 
\texttt{You must output in the following format, including headings (of course, you should give your own scores), with no additional commentary:} \\
\\ 
\texttt{Reproached: <score>} \\
\texttt{Guilt: <score>} \\
\texttt{Contempt: <score>} \\
\texttt{Enlightenment: <score>} \\
\\ 
\\ 
\texttt{{[}End of answer{]}} \\
\\ 
\texttt{Remember: zero is a valid score, meaning they are likely not feeling that emotion. You must score at least one emotion > 0.} \\
\\ 
\texttt{Your answer:} \\
            }
        }
    }
     \\ \\

    \theutterance \stepcounter{utterance}  
    & \multicolumn{4}{p{0.6\linewidth}}{
        \cellcolor[rgb]{0.561,0.737,0.561}{
            \makecell[{{p{\linewidth}}}]{
                \texttt{\tiny{[P1$\rangle$GM]}}
                \texttt{Reproached: 9} \\
\texttt{Guilt: 1} \\
\texttt{Contempt: 4} \\
\texttt{Enlightenment: 0} \\
            }
        }
    }
    & & \\ \\

    \theutterance \stepcounter{utterance}  
    & & & \multicolumn{2}{p{0.3\linewidth}}{
        \cellcolor[rgb]{0.9,0.9,0.9}{
            \makecell[{{p{\linewidth}}}]{
                \texttt{\tiny{[GM$|$GM]}}
                \texttt{Reproached: 9} \\
\texttt{Guilt: 1} \\
\texttt{Contempt: 4} \\
\texttt{Enlightenment: 0} \\
            }
        }
    }
    & & \\ \\

    \theutterance \stepcounter{utterance}  
    & & & \multicolumn{2}{p{0.3\linewidth}}{
        \cellcolor[rgb]{0.9,0.9,0.9}{
            \makecell[{{p{\linewidth}}}]{
                \texttt{\tiny{[GM$|$GM]}}
                \texttt{Reproached: 8} \\
\texttt{Guilt: 0} \\
\texttt{Contempt: 7} \\
\texttt{Enlightenment: 0} \\
            }
        }
    }
    & & \\ \\

    \theutterance \stepcounter{utterance}  
    & & & \multicolumn{2}{p{0.3\linewidth}}{
        \cellcolor[rgb]{0.9,0.9,0.9}{
            \makecell[{{p{\linewidth}}}]{
                \texttt{\tiny{[GM$|$GM]}}
                \texttt{game\_result = LOSE} \\
            }
        }
    }
    & & \\ \\

    \theutterance \stepcounter{utterance}  
    & & \multicolumn{4}{p{0.6\linewidth}}{
        \cellcolor[rgb]{0.95,0.95,0.95}{
            \makecell[{{p{\linewidth}}}]{
                \texttt{\tiny{[P1$\leftrightarrow$P]}}
                \texttt{Post{-}episode Reflexion (not sent to GM):} \\
\texttt{New plan: My previous attempts failed because I oscillated between overestimating and underestimating "Contempt" and "Guilt" without a consistent anchor to the text. The core mistake was treating Freddy's defensive rationalization ("Justice is a luxury") as a sign of either deep shame (Guilt) or genuine superiority (Contempt), when it is actually a mechanism to deflect the "Reproach." Freddy explicitly rejects the moral framework Izzy uses, calling her view a "luxury" he can't afford, which indicates he feels attacked (Reproached) but remains intellectually dismissive of her stance (Contempt). He does not show remorse (Guilt) because he believes his actions are necessary. Therefore, the correct strategy is to assign the highest score to "Reproached" (9{-}10) as the primary reaction to the "monster" accusation, a moderate score to "Contempt" (4{-}5) to reflect his dismissal of her idealism, and near{-}zero scores to "Guilt" (0{-}1) and "Enlightenment" (0). I will stop trying to find hidden guilt in his rationalizations and instead score based on his explicit rejection of her values.} \\
            }
        }
    }
    & \\ \\

\end{supertabular}
}

\end{document}
